%-------------------------------------------------------------------------------
%	WORK PLAN PAGE
%   Үзлэг, явц, шүүмжийн гарын үсэгтэй хуудсуудыг scan хийж signatures/huudas1..5.pdf нэрээр хадгалаад
%   доорх \includepdf мөрүүдийн %-ийг арилгана (text2sql-thesis/thesis/FILL_IN.md-тэй ижил).
%-------------------------------------------------------------------------------

\begin{titlepage}
\small
\noindent Батлав. \deptname ийн эрхлэгч:
\begin{flushright}
\makebox[6cm]{\dotfill} /\chairname/
\end{flushright}
Удирдагч:\makebox[6cm]{ \dotfill} /\supname/
\begin{center}
\vspace*{0.5cm}
\textbf{{\large \textsc{Дипломын төсөл гүйцэтгэх төлөвлөгөө}}}\\[0.5cm]
\end{center}
\noindent \textbf{Дипломын төслийн сэдэв:}\\
\textbf{Монгол}: «\ttitle» \\
\textbf{Англи}: «\ttitleng»\\

\noindent \textbf{Төслийн зорилго}: Монгол хэлээр (кирилл, латин галиг) асуусан асуултыг банк бус санхүүгийн байгууллагын
PostgreSQL өгөгдлийн сан дээр гүйцэтгэх SQL болгон хөрвүүлэх, нээлттэй эхийн загварт суурилсан системийг судалгааны
үр дүнд тулгуурлан бүтээж, хяналттай туршилтаар үнэлэх.\\

\noindent \textbf{Гүйцэтгэх оюутны овог нэр}:\makebox[4cm]{ } \shortname /\studentcode/ \\
\textbf{Холбоо барих утас}: \makebox[7cm]{ }\phonenum \\
\noindent
	\begin{tabularx}{\textwidth}{|>{\centering\arraybackslash}p{0.6cm}|X|>{\centering\arraybackslash}p{1.3cm}|>{\centering\arraybackslash}p{1.8cm}|}
	\hline
	\textbf{№} & \textbf{Ажлын бүлэг, хэсгийн нэр} & \textbf{Эзлэх хувь} & \textbf{Дуусах хугацаа} \\ \hline
	\multicolumn{4}{|l|}{\textbf{Бүлэг №1. Text2SQL-ийн онол ба судалгааны тойм}} \\ \hline
	1 & 1.1 Text2SQL асуудал ба хөгжил\newline
	    1.2 Том хэлний загварт суурилсан архитектур\newline
	    1.3 Хайлтаар баяжуулсан үүсгэлт (RAG)\newline
	    1.4 Fine-tuning: SFT ба бататгах сургалт\newline
	    1.5 Agentic Text2SQL\newline
	    1.6 Бага нөөцтэй хэл ба монгол хэлний сорилт\newline
	    1.7 Үнэлгээ ба аюулгүй байдал\newline
	    1.8 Бүлгийн дүгнэлт & 30\% & end5896 \\ \hline
	\multicolumn{4}{|l|}{\textbf{Бүлэг №2. Системийн зохиомж ба хэрэгжүүлэлт}} \\ \hline
	2 & 2.1 Шаардлага ба ерөнхий архитектур\newline
	    2.2 Өгөгдлийн сан ба text2sql давхарга\newline
	    2.3 Мэдлэгийн давхарга ба хайлт\newline
	    2.4 Загварын протокол ба SQL хамгаалалт\newline
	    2.5 Agentic урсгал\newline
	    2.6 Fine-tuning-ийн дамжлага\newline
	    2.7 Үйлчилгээ ба аюулгүй байдал\newline
	    2.8 Бүлгийн дүгнэлт & 30\% & end5896 \\ \hline
	\multicolumn{4}{|l|}{\textbf{Бүлэг №3. Туршилт ба үр дүн}} \\ \hline
	3 & 3.1 Үнэлгээний асуултын сан\newline
	    3.2 Үнэлгээний арга зүй\newline
	    3.3 Туршилтын төлөвлөгөө\newline
	    3.4 Системийн баталгаажуулалт\newline
	    3.5 Туршилтын үр дүн\newline
	    3.6 text2sql-thesis-тэй харьцуулалт\newline
	    3.7 Хязгаарлалт\newline
	    3.8 Бүлгийн дүгнэлт & 40\% & end5896 \\ \hline
	\multicolumn{4}{|l|}{\textbf{Ерөнхий дүгнэлт ба хавсралт}} \\ \hline
\end{tabularx}

\vspace{0.5cm}
Төлөвлөгөөг боловсруулсан оюутан: \makebox[3cm]{\dotfill} /\shortname/
\end{titlepage}

% Гарын үсэгтэй scan хуудсууд (байгаа бол %-ийг арилга):
%\includepdf[pages=-]{signatures/huudas1.pdf}
%\includepdf[pages=-]{signatures/huudas2.pdf}
%\includepdf[pages=-]{signatures/huudas3.pdf}
%\includepdf[pages=-]{signatures/huudas4.pdf}
%\includepdf[pages=-]{signatures/huudas5.pdf}
